% ============================================================================
% CAPITULO 7: MATEMATICA, TEORIA DOS NUMEROS, COMBINATORIA E FFT/NTT
% ============================================================================
\section{Matemática, Teoria dos Números, Combinatória e FFT/NTT}

% ----------------------------------------------------------------------------
% SUBTEMA 7.1: MANIPULACAO DE BITS, RECURSAO BINARIA E DIVISORES O(SQRT(N))
% ----------------------------------------------------------------------------
\subsection{Manipulação de Bits, Recursão Binária e Divisores em $\Ocomp{\sqrt{N}}$}

\subsubsection{Rastreamento Reverso de Bits em Dobras de Papel}
\algobox
  {Dada uma tira de $2^N$ seções ($N \le 60$), encontrar a sequência de dobras (\texttt{'L'} ou \texttt{'R'}) que leva a seção inicial $P \in [1, 2^N]$ para a altura final $H \in [1, 2^N]$.}
  {Trabalhando em 0-based ($p = P - 1, h = H - 1 \in [0, 2^N - 1]$):
  $\bullet$ Olhando do final para o início (desdobrando a pilha): no último passo restam $2^N$ camadas na altura $h$. Se $h$ está na metade superior da pilha atual ($h \ge 2^{i-1}$), essa camada veio da metade que foi \textbf{dobrada por cima} (e sua ordem vertical foi invertida: $h_{\text{ant}} = 2^i - 1 - h$). Se $h < 2^{i-1}$, ela estava na metade que ficou na base (sem inverter: $h_{\text{ant}} = h$).
  $\bullet$ Ao mesmo tempo, na tira original, no passo $k \in [1, N]$ (onde a tira tem largura $2^{N - k + 1}$), sabemos em qual metade (esquerda ou direita) a seção $p$ está olhando o bit $(N - k)$ de $p$ (mas toda vez que uma dobra inverte a nossa metade, o índice horizontal relativo dentro da metade também espelha: $p \gets 2^{N - k} - 1 - p$)!
  Logo, determinamos primeiro a sequência de bits \texttt{flipped[k]} ($k = 1 \dots N$) que diz se na $k$-ésima dobra a nossa camada subiu para o topo (foi invertida) ou ficou na base, inspecionando os bits de $h = H - 1$ de cima para baixo! Depois simulamos $p = P - 1$ da dobra $1$ até $N$: se $p$ está na metade esquerda e precisa ser invertido (\texttt{flipped[k] == 1}), dobramos \texttt{'L'}; se está na direita e precisa ser invertido, dobramos \texttt{'R'}!}
  {Usar \texttt{1ULL << b} (64 bits sem sinal) pois $N = 60$ atinge $2^{60}$.}
  {Tempo: $\Ocomp{N}$ $\mid$ Espaço: $\Ocomp{N}$}

\begin{lstlisting}
string solve_city_folding(int N, ull P, ull H) {
  ull p = P - 1, h = H - 1;
  /* 1. Descobre de tras para frente (passo N ate 1) se na dobra k nossa camada foi invertida */
  vector<bool> flip_step(N, false);
  for (int k = N - 1; k >= 0; --k) {
    ull half_h = 1ULL << k;
    if (h >= half_h) {
      flip_step[N - 1 - k] = true; /* Subiu para a metade superior (inverteu!) */
      h = (2 * half_h - 1) - h;
    }
  }
  /* 2. Simula do passo 0 ate N-1 acompanhando a posicao horizontal p */
  string ans(N, ' ');
  rep(k, 0, N) {
    ull half_w = 1ULL << (N - 1 - k);
    bool in_left = (p < half_w);
    bool need_flip = flip_step[k];
    if (need_flip) {
      ans[k] = in_left ? 'L' : 'R';
      p = in_left ? (half_w - 1 - p) : (2 * half_w - 1 - p);
    } else {
      ans[k] = in_left ? 'R' : 'L';
      if (!in_left) p -= half_w;
    }
  }
  return ans;
}
\end{lstlisting}

\subsubsection{Bases Palindrômicas e Enumeração de Divisores em $\Ocomp{\sqrt{N}}$}
\algobox
  {Dado $N \le 10^{12}$, encontrar todas as bases $b \in [2, N]$ nas quais a representação de $N$ na base $b$ é um palíndromo.}
  {Dividimos a busca no limiar $\sqrt{N} \le 10^6$:
  $\bullet$ \textbf{Caso 1 ($2 \le b \le \sqrt{N}$):} Existem no máximo $10^6$ bases; para cada uma, extraímos os dígitos de $N$ na base $b$ em $\Ocomp{\log_b N}$ e testamos se o vetor de dígitos é palindrômico.
  $\bullet$ \textbf{Caso 2 ($b > \sqrt{N}$):} Como $b > \sqrt{N}$ e $b \le N$, o número $N$ escrito na base $b$ possui \textbf{exatamente 2 dígitos}: $N = (d_1 d_0)_b = d_1 \cdot b + d_0$, com $1 \le d_1 < b$ e $0 \le d_0 < b$. Para ser um palíndromo de 2 dígitos, devemos ter $d_1 = d_0 = d$, ou seja:
  $\displaystyle N = d \cdot b + d = d(b + 1) \implies (b + 1) = \frac{N}{d} \implies b = \frac{N}{d} - 1$
  onde $d$ é um divisor de $N$ satisfazendo $1 \le d < b = \frac{N}{d} - 1 \iff d(d + 1) < N$!
  Assim, basta enumerar os divisores $d \mid N$ com $d \le \sqrt{N}$ e adicionar $b = \frac{N}{d} - 1$ sempre que $b > \sqrt{N}$ (ou $d < b$)!}
  {Na base $b = N$, $N = (10)_N$, que nunca é palíndromo (termina em 0). Para $N = 2$, não há nenhuma base válida (retornar \texttt{*}).}
  {Tempo: $\Ocomp{\sqrt{N} \log N}$ $\mid$ Espaço: $\Ocomp{\sigma_0(N)}$}

\begin{lstlisting}
vll palindromic_bases(ll N) {
  vll ans;
  auto is_pal_base = [](ll n, ll b) -> bool {
    vll digs;
    while (n > 0) { digs.pb(n % b); n /= b; }
    int m = sz(digs);
    rep(i, 0, m / 2) if (digs[i] != digs[m - 1 - i]) return false;
    return true;
  };
  ll sq = 1;
  while ((sq + 1) * (sq + 1) <= N) ++sq;

  /* 1. Testa diretamente todas as bases b in [2, sqrt(N)] */
  for (ll b = 2; b <= sq; ++b) {
    if (is_pal_base(N, b)) ans.pb(b);
  }
  /* 2. Para b > sqrt(N), N na base b tem 2 digitos iguais: N = d * b + d => b = N/d - 1 */
  vll large_bases;
  for (ll d = 1; d <= sq; ++d) {
    if (N % d == 0) {
      ll b = (N / d) - 1;
      if (b > sq && d < b) large_bases.pb(b);
    }
  }
  reverse(all(large_bases));
  ans.insert(ans.end(), all(large_bases));
  return ans;
}
\end{lstlisting}

% ----------------------------------------------------------------------------
% SUBTEMA 7.2: COMBINATORIA, INVERSOES, POLINOMIOS E CICLOS DE PERMUTACAO
% ----------------------------------------------------------------------------
\subsection{Combinatória Modular, Inversões, Polinômios e Ciclos}

\subsubsection{Aritmética Modular (\texttt{Mint}), Euclides Estendido, CRT e Binomiais}
\algobox
  {(1) Struct \texttt{Mint} para operações $\bmod M$; (2) Euclides Estendido ($ax + by = \gcd(a,b)$) e Teorema Chinês do Resto (CRT) para módulos não-coprimos; (3) Fatoriais e combinações $\binom{n}{k}$ em $\Ocomp{1}$.}
  {No CRT geral para $x \equiv r_1 \pmod{m_1}$ e $x \equiv r_2 \pmod{m_2}$, existe solução sse $g = \gcd(m_1, m_2)$ divide $(r_2 - r_1)$.}
  {Usar \texttt{\_\_int128} na multiplicação do CRT para evitar overflow quando $\text{lcm}(m_1, m_2) \le 10^{18}$.}
  {Euclides / CRT: $\Ocomp{\log(\min(a,b))}$ $\mid$ Binomial: $\Ocomp{1}$ após pré-cálculo $\Ocomp{N}$}

\begin{lstlisting}
template <ll MOD = 1'000'000'007>
struct Mint {
  ll v;
  Mint(ll _v = 0) : v((_v % MOD + MOD) % MOD) {}
  Mint& operator+=(const Mint& o) { if ((v += o.v) >= MOD) v -= MOD; return *this; }
  Mint& operator-=(const Mint& o) { if ((v -= o.v) < 0) v += MOD; return *this; }
  Mint& operator*=(const Mint& o) { v = (v * o.v) % MOD; return *this; }
  Mint pow(ll e) const {
    Mint r = 1, b = *this;
    for (; e > 0; e >>= 1, b *= b) if (e & 1) r *= b;
    return r;
  }
  Mint inv() const { return pow(MOD - 2); }
  Mint& operator/=(const Mint& o) { return *this *= o.inv(); }
  friend Mint operator+(Mint a, const Mint& b) { return a += b; }
  friend Mint operator-(Mint a, const Mint& b) { return a -= b; }
  friend Mint operator*(Mint a, const Mint& b) { return a *= b; }
  friend Mint operator/(Mint a, const Mint& b) { return a /= b; }
};

struct Comb {
  int n;
  vector<Mint<>> fact, invf;
  Comb(int n) : n(n), fact(n + 1, 1), invf(n + 1, 1) {
    rep(i, 1, n + 1) fact[i] = fact[i - 1] * i;
    invf[n] = fact[n].inv();
    for (int i = n - 1; i >= 0; --i) invf[i] = invf[i + 1] * (i + 1);
  }
  Mint<> C(int a, int b) const {
    if (b < 0 || b > a) return 0;
    return fact[a] * invf[b] * invf[a - b];
  }
};

/* Euclides Estendido: retorna g = gcd(a, b) e preenche x, y tais que a*x + b*y = g */
ll extgcd(ll a, ll b, ll& x, ll& y) {
  if (!b) { x = 1; y = 0; return a; }
  ll x1, y1, g = extgcd(b, a % b, x1, y1);
  x = y1; y = x1 - y1 * (a / b);
  return g;
}

/* CRT Geral (modulos nao necessariamente coprimos): retorna {rem, mod} ou {-1, -1} */
pll crt_merge(ll r1, ll m1, ll r2, ll m2) {
  ll x, y, g = extgcd(m1, m2, x, y);
  if ((r2 - r1) % g != 0) return {-1, -1};
  ll lcm = (m1 / g) * m2;
  ll mult = (i128)((r2 - r1) / g) * x % (m2 / g);
  ll r = (r1 + (i128)mult * m1) % lcm;
  if (r < 0) r += lcm;
  return {r, lcm};
}
\end{lstlisting}

\subsubsection{Expansão Polinomial $\prod (x - r_k)$ em $\Ocomp{D^2}$ e Inversões em $S^N$}
\algobox
  {(1) Para construir polinômio por sinais, construir o polinômio de menor grau com coeficientes inteiros e coeficiente líder $\pm 1$ tal que $P(2i) > 0$ se $S_i = \text{'H'}$ e $P(2i) < 0$ se $S_i = \text{'A'}$; (2) Para inversões em string repetida, contar inversões na string $S$ repetida $N \le 10^{12}$ vezes $\bmod 10^9+7$.}
  {\textbf{1. Polinômio por Trocas de Sinal:} O polinômio deve mudar de sinal entre $2i$ e $2(i+1)$ sse $S[i] \neq S[i+1]$! Colocamos uma raiz inteira ímpar $r_k = 2i + 1$ exatamente em cada transição $S[i] \neq S[i+1]$, obtendo o menor grau possível $D$ (número de trocas de sinal). Multiplicamos $\prod_{k=1}^D (x - r_k)$ em $\Ocomp{D^2}$ e ajustamos o sinal global de todos os coeficientes olhando se para $x = 2N$ (maior que todas as raízes, onde $\prod(x - r_k) > 0$) o último caractere $S[N]$ é \texttt{'H'} ($+1$) ou \texttt{'A'} ($-1$)!
  \textbf{2. Inversões em $S^N$:} Uma inversão em $S^N$ ou ocorre dentro da mesma cópia de $S$ ($N \cdot \text{inv}_{\text{in}}(S)$), ou ocorre entre uma letra $S[i]$ na cópia $a$ e uma letra menor $S[j] < S[i]$ na cópia $b > a$ ($\binom{N}{2} \cdot \sum_{c_1 > c_2} \text{cnt}[c_1] \cdot \text{cnt}[c_2]$)!}
  {No polinômio, os coeficientes cabem em \texttt{signed long long} ($< 2^{63}$ garantido pelo enunciado). Nas inversões em $S^N$, reduzir $N \pmod{10^9+7}$ antes de calcular $\binom{N}{2} = N(N-1)/2 \pmod{10^9+7}$.}
  {Polinômio: $\Ocomp{D^2}$ $\mid$ Inversões em $S^N$: $\Ocomp{|S|}$}

\begin{lstlisting}
/* 1. Polinomio por trocas de sinal: Retorna {grau D, coeficientes de x^D ate x^0} */
pair<int, vll> build_alien_polynomial(const string& S) {
  int n = sz(S);
  vll roots;
  rep(i, 0, n - 1) if (S[i] != S[i + 1]) roots.pb(2LL * (i + 1) + 1);
  int D = sz(roots);
  /* poly[k] = coeficiente de x^k (0 <= k <= D) */
  vll poly(D + 1, 0);
  poly[0] = 1;
  int cur_deg = 0;
  for (ll r : roots) {
    ++cur_deg;
    for (int k = cur_deg; k >= 1; --k) poly[k] = poly[k - 1] - r * poly[k];
    poly[0] *= -r;
  }
  /* Para x = 2*n (acima de todas as raizes), prod(x - r) > 0.
   * Se o ultimo cidadao S[n-1] for 'A', multiplicamos todos os coeficientes por -1! */
  if (S.back() == 'A') {
    for (ll& c : poly) c = -c;
  }
  reverse(all(poly)); /* Ordem decrescente de grau: x^D, x^(D-1), ..., x^0 */
  return {D, poly};
}

/* 2. Inversoes na string S repetida N vezes mod 1e9+7 */
ll count_repeated_string_inversions(const string& S, ll N) {
  constexpr ll MOD = 1'000'000'007, INV2 = (MOD + 1) / 2;
  vll cnt(26, 0);
  ll inv_inside = 0;
  for (char ch : S) {
    int c = ch - 'a';
    rep(bigger, c + 1, 26) inv_inside = (inv_inside + cnt[bigger]) % MOD;
    cnt[c]++;
  }
  ll pairs_diff_copies = 0;
  rep(c1, 0, 26) rep(c2, 0, c1) {
    pairs_diff_copies = (pairs_diff_copies + cnt[c1] * cnt[c2]) % MOD;
  }
  ll n_mod = N % MOD;
  ll c2_N = n_mod * ((n_mod - 1 + MOD) % MOD) % MOD * INV2 % MOD;
  return (n_mod * inv_inside + c2_N * pairs_diff_copies) % MOD;
}
\end{lstlisting}

\subsubsection{Ciclos de Permutação e Quadrado Latino Invariante}
\algobox
  {Dada uma sequência de trocas (transposições) de linhas e de colunas sobre uma matriz $N \times N$, decidir se existe um Quadrado Latino que permanece idêntico a si mesmo após todas as trocas e construí-lo.}
  {Aplicar transposições de linhas equivale a permutar as linhas por uma permutação $\pi_R$, e as colunas por uma permutação $\pi_C$. O Quadrado Latino $M$ é invariante ($M_{\pi_R(i), \pi_C(j)} = M_{i, j}$ para todo $i, j$) sse $M$ é constante ao longo das órbitas conjuntas $(i, j) \mapsto (\pi_R(i), \pi_C(j))$.
  $\bullet$ \textbf{Teorema de Existência:} Decomponha $\pi_R$ e $\pi_C$ em ciclos disjuntos! Para que um bloco $C_R \times C_C$ (onde $C_R$ é um ciclo de linhas de tamanho $a$ e $C_C$ é um ciclo de colunas de tamanho $b$) possa fazer parte de um Quadrado Latino invariante sem repetir símbolos na mesma linha ou coluna, todas as órbitas precisam ser compatíveis; especificamente, um Quadrado Latino invariante sob $(\pi_R, \pi_C)$ existe \textbf{se e somente se $\pi_R$ e $\pi_C$ possuem exatamente a mesma partição de comprimentos de ciclos} (ou seja, os multiconjuntos dos tamanhos dos ciclos de $\pi_R$ e de $\pi_C$ são idênticos)!
  $\bullet$ \textbf{Construção Explícita em $\Ocomp{N^2}$:} Quando os tamanhos dos ciclos coincidem ($L_1, L_2, \dots, L_m$), pareamos o $p$-ésimo ciclo de linhas $R_p = (r_{p,0}, \dots, r_{p, L_p-1})$ com o $p$-ésimo ciclo de colunas $C_p = (c_{p,0}, \dots, c_{p, L_p-1})$. Mais ainda: todos os ciclos de tamanho $s$ têm o mesmo tamanho $L_p$ (de fato, se uma linha está num ciclo de tamanho $a$ e uma coluna num ciclo de tamanho $b$ com $a \neq b$, após $a$ passos a linha volta ao início mas a coluna não, o que forçaria o mesmo símbolo a se repetir na mesma linha em colunas distintas se $a < b$, logo \textbf{todos os ciclos de $\pi_R$ e $\pi_C$ precisam ter o MESMO tamanho $L$ divisor de $N$}, e ter a mesma quantidade $K = N / L$ de ciclos!).
  Portanto: um Quadrado Latino invariante existe sse \textbf{todos os ciclos de $\pi_R$ e todos os ciclos de $\pi_C$ possuem exatamente o mesmo comprimento $L$}! E quando todos os $K = N / L$ ciclos de linhas e $K$ ciclos de colunas têm tamanho $L$, para o bloco entre o $u$-ésimo ciclo de linhas ($0 \le u < K$) e o $v$-ésimo ciclo de colunas ($0 \le v < K$), na posição interna $(i, j) \in [0, L-1]^2$, atribuímos o símbolo:
  $\displaystyle \text{val} = \bigl((u + v) \bmod K\bigr) \cdot L + \bigl((j - i + L) \bmod L\bigr) + 1$
  Como $(j + 1) - (i + 1) \equiv j - i \pmod L$, avançar simultaneamente no ciclo de linhas ($i+1$) e no ciclo de colunas ($j+1$) preserva o valor, e cada linha e coluna contém todos os números de $1$ a $N$!}
  {Simular os swaps rastreando de onde veio cada linha/coluna (ou para onde vai): na invariância $M' = M$, tanto $\pi$ quanto $\pi^{-1}$ têm os mesmos comprimentos de ciclo e as mesmas órbitas se girarmos ambos no mesmo sentido!}
  {Tempo: $\Ocomp{T + N^2}$ $\mid$ Espaço: $\Ocomp{N^2}$}

\begin{lstlisting}
/* Quadrado Latino invariante sob permutacoes de linhas e colunas em O(T + N^2)
 * Retorna a matriz N x N (1-based values) ou {} se impossivel ("*") */
vector<vi> build_invariant_latin_square(int N, const vector<tuple<char, int, int>>& ops) {
  /* posR[i] = qual linha original esta atualmente na posicao i (0..N-1) */
  vi posR(N), posC(N);
  iota(all(posR), 0); iota(all(posC), 0);
  for (auto [type, u, v] : ops) {
    --u; --v;
    if (type == 'R') swap(posR[u], posR[v]);
    else             swap(posC[u], posC[v]);
  }
  /* A linha que comecou em posR[i] terminou na linha i: nextR[posR[i]] = i */
  vi nextR(N), nextC(N);
  rep(i, 0, N) { nextR[posR[i]] = i; nextC[posC[i]] = i; }

  auto get_cycles = [&](const vi& nxt) -> vector<vi> {
    vector<vi> cycs;
    vector<bool> vis(N, false);
    rep(i, 0, N) if (!vis[i]) {
      vi c;
      for (int cur = i; !vis[cur]; cur = nxt[cur]) {
        vis[cur] = true; c.pb(cur);
      }
      cycs.pb(move(c));
    }
    return cycs;
  };
  auto cycR = get_cycles(nextR), cycC = get_cycles(nextC);
  int L = sz(cycR[0]);
  for (const auto& c : cycR) if (sz(c) != L) return {};
  for (const auto& c : cycC) if (sz(c) != L) return {};

  int K = N / L; /* Quantidade de ciclos de tamanho L */
  vector<vi> ans(N, vi(N, 0));
  rep(u, 0, K) rep(v, 0, K) {
    int base = ((u + v) % K) * L;
    rep(i, 0, L) rep(j, 0, L) {
      int r = cycR[u][i], c = cycC[v][j];
      ans[r][c] = base + ((j - i + L) % L) + 1;
    }
  }
  return ans;
}
\end{lstlisting}

% ----------------------------------------------------------------------------
% SUBTEMA 7.3: FFT / NTT, ORDEM MULTIPLICATIVA, TOTIENTE E POLLARD'S RHO
% ----------------------------------------------------------------------------
\subsection{FFT / NTT, Ordem Multiplicativa, Totiente e Pollard's Rho}

\subsubsection{Fast Fourier Transform (FFT), NTT e Pareamento Extremo + Convolução}
\algobox
  {(1) Multiplicar dois polinômios em $\Ocomp{N \log N}$ sobre $\mathbb{C}$ (FFT para inteiros sem módulo) ou sobre $\ZZ_{998244353}$ (NTT exato); (2) Para produtos escalares extremos, dados dois vetores $F$ e $C$ de tamanho $N \le 10^5$ com valores em $[-10^4, 10^4]$, encontrar para todo $k \in [1, N]$ o menor e o maior produto escalar de tamanho $k$ escolhendo $k$ elementos distintos de $F$ e $k$ de $C$.}
  {\textbf{Pareamento Extremo + Convolução:}
  $\bullet$ \textbf{Para o Máximo:} Ordenamos $F$ e $C$. Os elementos positivos de $F$ querem multiplicar os positivos de $C$ (do maior para o maior: $F_{\text{pos}} \times C_{\text{pos}}$), e os negativos de $F$ querem multiplicar os negativos de $C$ (do mais negativo para o mais negativo: $F_{\text{neg}} \times C_{\text{neg}}$)! Seja $P_{\text{same}}$ a lista desses produtos positivos ordenados do maior para o menor: para $k \le |P_{\text{same}}|$, basta somar os $k$ maiores produtos de $P_{\text{same}}$!
  E quando $k > |P_{\text{same}}|$? Já esgotamos todos os pares de mesmo sinal! Os elementos que sobraram em $F$ têm todos o mesmo sinal (digamos, positivos em $F$) e os que sobraram em $C$ têm todos o sinal oposto (negativos em $C$)! Para minimizar a perda ao escolher $m = k - |P_{\text{same}}|$ pares de sinais opostos, no conjunto de pares de mesmo sinal nunca devemos desfazer um par (já que trocar mesmo sinal por sinais opostos só diminui a soma), e dentre os elementos restantes de sinais opostos (ordenados por valor absoluto crescente: $A_0 \le A_1 \dots$ e $B_0 \le B_1 \dots$), para qualquer $m$ fixo, pela desigualdade de rearranjo devemos escolher os $m$ menores valores absolutos $A_0 \dots A_{m-1}$ e $B_0 \dots B_{m-1}$ pareados em \textbf{ordem reversa}: $\sum_{i=0}^{m-1} A_i \cdot B_{m - 1 - i}$!
  Veja a fórmula $\sum_{i=0}^{m-1} A_i B_{m-1-i}$: isso é \textbf{exatamente o coeficiente $x^{m-1}$ da Convolução Polinomial (FFT)} dos polinômios $A(x) = \sum A_i x^i$ e $B(x) = \sum B_i x^i$!
  $\bullet$ \textbf{Para o Mínimo:} Basta negar todos os elementos de $C$ ($C_i \gets -C_i$), rodar exatamente a mesma função de Máximo e negar o resultado!}
  {Usar \texttt{double} ou \texttt{long double} na FFT (como $|F_i C_j| \le 10^8$ e $N \le 10^5$, o coeficiente máximo da convolução é $10^{13}$, que cabe com folga nos 53 bits de mantissa de \texttt{double} $\approx 9 \cdot 10^{15}$, arredondando com \texttt{llround}).}
  {Tempo: $\Ocomp{N \log N}$ $\mid$ Espaço: $\Ocomp{N}$}

\begin{lstlisting}
/* 1. FFT sobre Complex<double> e Convolucao Inteira Exata (ate ~10^15) */
using cd = complex<double>;

void fft(vector<cd>& a, bool invert) {
  int n = sz(a);
  for (int i = 1, j = 0; i < n; ++i) {
    int bit = n >> 1;
    for (; j & bit; bit >>= 1) j ^= bit;
    j ^= bit;
    if (i < j) swap(a[i], a[j]);
  }
  for (int len = 2; len <= n; len <<= 1) {
    double ang = 2 * acos(-1.0) / len * (invert ? -1 : 1);
    cd wlen(cos(ang), sin(ang));
    for (int i = 0; i < n; i += len) {
      cd w(1);
      rep(j, 0, len / 2) {
        cd u = a[i + j], v = a[i + j + len / 2] * w;
        a[i + j] = u + v;
        a[i + j + len / 2] = u - v;
        w *= wlen;
      }
    }
  }
  if (invert) for (cd& x : a) x /= n;
}

vll poly_multiply_fft(const vll& a, const vll& b) {
  if (a.empty() || b.empty()) return {};
  int s = sz(a) + sz(b) - 1, n = 1;
  while (n < s) n <<= 1;
  vector<cd> fa(n, 0), fb(n, 0);
  rep(i, 0, sz(a)) fa[i] = (double)a[i];
  rep(i, 0, sz(b)) fb[i] = (double)b[i];
  fft(fa, false); fft(fb, false);
  rep(i, 0, n) fa[i] *= fb[i];
  fft(fa, true);
  vll res(s);
  rep(i, 0, s) res[i] = llround(fa[i].real());
  return res;
}

/* 2. Number Theoretic Transform (NTT) exato mod 998244353 (Raiz primitiva G = 3) */
void ntt(vll& a, bool invert) {
  static constexpr ll MOD = 998244353, G = 3;
  auto modpow = [](ll b, ll e) {
    ll r = 1;
    for (; e > 0; e >>= 1, b = b * b % MOD) if (e & 1) r = r * b % MOD;
    return r;
  };
  int n = sz(a);
  for (int i = 1, j = 0; i < n; ++i) {
    int bit = n >> 1;
    for (; j & bit; bit >>= 1) j ^= bit;
    j ^= bit;
    if (i < j) swap(a[i], a[j]);
  }
  for (int len = 2; len <= n; len <<= 1) {
    ll wlen = modpow(G, (MOD - 1) / len);
    if (invert) wlen = modpow(wlen, MOD - 2);
    for (int i = 0; i < n; i += len) {
      ll w = 1;
      rep(j, 0, len / 2) {
        ll u = a[i + j], v = (a[i + j + len / 2] * w) % MOD;
        a[i + j] = (u + v < MOD ? u + v : u + v - MOD);
        a[i + j + len / 2] = (u - v >= 0 ? u - v : u - v + MOD);
        w = (w * wlen) % MOD;
      }
    }
  }
  if (invert) {
    ll n_inv = modpow(n, MOD - 2);
    for (ll& x : a) x = (x * n_inv) % MOD;
  }
}

/* 3. Menor e maior produto escalar de tamanho k para k = 1..N em O(N log N) */
vll max_subset_dot_products(vll F, vll C) {
  int n = sz(F);
  vll F_pos, F_neg, C_pos, C_neg;
  for (ll x : F) (x >= 0 ? F_pos : F_neg).pb(llabs(x));
  for (ll x : C) (x >= 0 ? C_pos : C_neg).pb(llabs(x));
  sort(all(F_pos), greater<ll>()); sort(all(C_pos), greater<ll>());
  sort(all(F_neg), greater<ll>()); sort(all(C_neg), greater<ll>());

  int kp = min(sz(F_pos), sz(C_pos)), kn = min(sz(F_neg), sz(C_neg));
  vll same_prods;
  same_prods.reserve(kp + kn);
  rep(i, 0, kp) same_prods.pb(F_pos[i] * C_pos[i]);
  rep(i, 0, kn) same_prods.pb(F_neg[i] * C_neg[i]);
  sort(all(same_prods), greater<ll>());

  /* Elementos restantes de sinais opostos (ordenados por valor absoluto CRESCENTE) */
  vll remA, remB;
  rep(i, kp, sz(F_pos)) remA.pb(F_pos[i]);
  rep(i, kn, sz(F_neg)) remA.pb(F_neg[i]);
  rep(i, kp, sz(C_pos)) remB.pb(C_pos[i]);
  rep(i, kn, sz(C_neg)) remB.pb(C_neg[i]);
  sort(all(remA)); sort(all(remB));

  vll conv = poly_multiply_fft(remA, remB);
  vll ans(n, 0);
  ll sum_same = 0;
  int S = sz(same_prods);
  rep(k, 1, n + 1) {
    if (k <= S) {
      sum_same += same_prods[k - 1];
      ans[k - 1] = sum_same;
    } else {
      int m = k - S; /* Precisamos de 'm' pares de sinais opostos (soma minima = conv[m - 1]) */
      ans[k - 1] = sum_same - conv[m - 1];
    }
  }
  return ans;
}

vector<pll> solve_because_art(const vll& F, vll C) {
  vll mx = max_subset_dot_products(F, C);
  for (ll& x : C) x = -x;
  vll mn = max_subset_dot_products(F, C);
  vector<pll> res(sz(F));
  rep(i, 0, sz(F)) res[i] = {-mn[i], mx[i]};
  return res;
}
\end{lstlisting}

\subsubsection{Crivo Linear, Função Totiente ($\phi$), Ordem Multiplicativa e Dízimas}
\algobox
  {(1) Crivo Linear até $M$ pré-calculando menor fator primo \texttt{spf[x]} e $\phi(x)$ em $\Ocomp{M}$; (2) Calcular a \textbf{Ordem Multiplicativa} $\text{ord}_m(10)$ (menor $d \ge 1$ tal que $10^d \equiv 1 \pmod m$ para $\gcd(m, 10) = 1$) testando divisores de $\phi(m)$; (3) Para dígitos de dízimas $A_i/A_j$, calcular o valor esperado do número de dígitos da representação decimal exata/periódica de $A_i / A_j$ ($\bmod 998244353$).}
  {\textbf{Análise Exata do Número de Dígitos de $\frac{A_i}{A_j}$ nas dízimas:}
  Reduza a fração $\frac{A_i}{A_j} = \frac{p}{q}$ dividindo por $g = \gcd(A_i, A_j)$.
  1. \textbf{Parte Inteira:} A parte inteira é $\lfloor p / q \rfloor$. Quantos dígitos ela usa? Se $\lfloor p / q \rfloor == 0$, usa $1$ dígito (\texttt{"0"}). Se $\lfloor p / q \rfloor \ge 1$, usa $\lfloor \log_{10}(p / q) \rfloor + 1$ dígitos!
  2. \textbf{Parte Fracionária (Pré-período + Período):} Escreva o denominador reduzido como $q = 2^a \cdot 5^b \cdot m$, onde $\gcd(m, 10) = 1$:
     $\bullet$ O número de dígitos decimais \textbf{não-periódicos} após a vírgula é exatamente $\max(a, b)$!
     $\bullet$ O número de dígitos do \textbf{período} da dízima é $0$ se $m = 1$, e é exatamente a \textbf{Ordem Multiplicativa $\text{ord}_m(10)$} se $m > 1$!
  Logo, para $\frac{p}{q}$ irredutível com $q = 2^a 5^b m$:
  $\displaystyle \text{digits}(p, q) = \text{int\_digits}(\lfloor p / q \rfloor) + \max(a, b) + (m > 1 ? \text{ord}_m(10) : 0)$
  \textbf{Agregação em $\Ocomp{M \log M}$ ($M = \max A_i \le 10^5$):}
  Seja $f[x]$ a frequência do valor $x$ no array $A$.
  $\bullet$ A soma de $\text{int\_digits}(\lfloor A_i / A_j \rfloor)$ para todos os pares $(i, j)$ é calculada em $\Ocomp{M}$ somando 1 para todos os $N^2$ pares, mais 1 para cada par com $A_i \ge 10 A_j$, mais 1 para $A_i \ge 100 A_j$, $\dots$, usando soma de prefixos de $f$!
  $\bullet$ O termo $\max(a, b) + \text{período}(m)$ depende \textbf{apenas do denominador reduzido} $q = \frac{A_j}{\gcd(A_i, A_j)}$! Pré-calculamos $C(q) = \max(a, b) + (m > 1 ? \text{ord}_m(10) : 0)$ para todo $q \in [1, M]$. Para somar $C(A_j / \gcd(A_i, A_j))$ sobre todos os pares $(i, j)$, usamos Inversão de Möbius / Convolução de GCD sobre os múltiplos de $g$: para cada $g \in [1, M]$, contamos quantos $A_i$ são múltiplos de $g$ e somamos pelos múltiplos $A_j = k \cdot g$ em $\Ocomp{M \log M}$!}
  {Na soma por GCD, um par $(x = u \cdot g, y = v \cdot g)$ com $\gcd(u, v) = 1$ tem denominador reduzido $q = v$! Logo, podemos pré-calcular para cada $v \in [1, M]$ o peso $W(v) = C(v)$ e usar Möbius $\mu(d)$ para filtrar $\gcd(u, v) = 1$: $\sum_{\gcd(u,v)=1} cnt_g(u) cnt_g(v) C(v) = \sum_{d} \mu(d) \Bigl(\sum_{d \mid u} cnt_g(u)\Bigr) \Bigl(\sum_{d \mid v} cnt_g(v) C(v)\Bigr)$ em $\Ocomp{M \log^2 M}$ ($\approx 1.5 \cdot 10^6$ operações, $< 0.02\text{s}$)!}
  {Tempo: $\Ocomp{N + M \log^2 M}$ ($M = 10^5$) $\mid$ Espaço: $\Ocomp{M}$}

\begin{lstlisting}
struct SieveAndOrder {
  int M;
  vi spf, phi, mu, primes;

  SieveAndOrder(int M) : M(M), spf(M + 1, 0), phi(M + 1, 0), mu(M + 1, 0) {
    phi[1] = 1; mu[1] = 1;
    rep(i, 2, M + 1) {
      if (spf[i] == 0) {
        spf[i] = i; phi[i] = i - 1; mu[i] = -1;
        primes.pb(i);
      }
      for (int p : primes) {
        if (p > spf[i] || i * p > M) break;
        spf[i * p] = p;
        if (p == spf[i]) { phi[i * p] = phi[i] * p; mu[i * p] = 0; }
        else             { phi[i * p] = phi[i] * (p - 1); mu[i * p] = -mu[i]; }
      }
    }
  }
  ll modpow(ll b, ll e, ll mod) const {
    ll r = 1 % mod;
    for (; e > 0; e >>= 1, b = b * b % mod) if (e & 1) r = r * b % mod;
    return r;
  }
  /* Ordem multiplicativa de 'a' modulo 'm' (requer gcd(a, m) == 1 e m <= M) */
  int mult_order(ll a, int m) const {
    if (m <= 1) return 0;
    int ord = phi[m], tmp = ord;
    while (tmp > 1) {
      int p = spf[tmp];
      while (tmp % p == 0) tmp /= p;
      while (ord % p == 0 && modpow(a, ord / p, m) == 1) ord /= p;
    }
    return ord;
  }
};

/* Valor esperado de digitos na expansao decimal de A_i / A_j em O(N + M log^2 M) */
ll solve_displaying_decimals(const vi& A) {
  constexpr ll MOD = 998244353;
  int N = sz(A), M = *max_element(all(A));
  SieveAndOrder sv(M);

  /* 1. Pre-calcula frac_digits[q] = max(a, b) + (m > 1 ? ord_m(10) : 0) para q = 2^a * 5^b * m */
  vll frac_dig(M + 1, 0);
  rep(q, 1, M + 1) {
    int m = q, a = 0, b = 0;
    while (m % 2 == 0) { m /= 2; ++a; }
    while (m % 5 == 0) { m /= 5; ++b; }
    frac_dig[q] = max(a, b) + (m > 1 ? sv.mult_order(10, m) : 0);
  }

  vll freq(M + 1, 0), pref(M + 1, 0);
  for (int x : A) freq[x]++;
  rep(i, 1, M + 1) pref[i] = pref[i - 1] + freq[i];

  /* 2. Soma dos digitos da parte inteira floor(A_i / A_j): */
  ll total_digits = (1LL * N * N) % MOD;
  rep(y, 1, M + 1) if (freq[y] > 0) {
    for (ll p10 = 10; y * p10 <= M; p10 *= 10) {
      ll cnt_x = pref[M] - pref[y * p10 - 1];
      total_digits = (total_digits + freq[y] * cnt_x) % MOD;
    }
  }

  /* 3. Soma de frac_dig[ A_j / gcd(A_i, A_j) ]: */
  vll mult_cnt(M + 1, 0);
  rep(l, 1, M + 1) {
    for (int m = l; m <= M; m += l) mult_cnt[l] += freq[m];
  }
  rep(k, 1, M + 1) if (sv.mu[k] != 0) {
    for (int vp = 1; vp * k <= M; ++vp) {
      ll term_v = frac_dig[vp * k];
      if (!term_v) continue;
      int step = vp * k; /* Precisamos somar sobre g = 1 .. M / (vp * k), onde l = g * k e y = g * k * vp */
      ll ways_y_x = 0;
      for (int g = 1; g * step <= M; ++g) {
        int l = g * k, y = g * step;
        if (freq[y] && mult_cnt[l]) ways_y_x += mult_cnt[l] * freq[y];
      }
      ways_y_x %= MOD;
      if (sv.mu[k] == 1) total_digits = (total_digits + ways_y_x * term_v) % MOD;
      else               total_digits = (total_digits - ways_y_x * term_v % MOD + MOD) % MOD;
    }
  }
  ll inv_N2 = sv.modpow(1LL * N * N % MOD, MOD - 2, MOD);
  return total_digits * inv_N2 % MOD;
}
\end{lstlisting}

\subsubsection{Miller-Rabin 64-bit, Pollard's Rho e Fatoração Difusa}
\algobox
  {(1) Testar primalidade determinística em 64 bits ($< 2^{64}$) em $\Ocomp{\log^3 N}$ (Miller-Rabin com as 7 bases de Sinclair); (2) Fatorar qualquer inteiro de 64 bits em $\Ocomp{N^{1/4}}$ via \textbf{Pollard's Rho}; (3) Para fatoração aproximada de inteiros gigantes, dado $X \le 10^{1000}$, encontrar a fatoração prima de um inteiro $Y$ com erro relativo $\frac{|X - Y|}{X} \le 10^{-9}$ cujos fatores primos sejam todos $\le 10^{18}$.}
  {\textbf{Fatoração Aproximada por Prefixo:}
  $\bullet$ Se $|S_X| \le 18$ ($X < 10^{18}$): basta fatorar $X$ exatamente com Pollard's Rho (erro relativo 0, todos os fatores primos $\le X \le 10^{18}$)!
  $\bullet$ Se $|S_X| > 18$: extraímos o prefixo $P$ formado pelos primeiros $18$ dígitos de $X$, e seja $E = |S_X| - 18$ a quantidade de dígitos restantes no sufixo. Qualquer número da forma $Y = P' \cdot 10^E = P' \cdot 2^E \cdot 5^E$ onde $|P' - P| \le 10^7$ satisfaz:
  $\displaystyle \frac{|X - Y|}{X} \le \frac{(|P' - P| + 1) \cdot 10^E}{P \cdot 10^E} \le \frac{10^7 + 1}{10^{17}} \approx 10^{-10} < 10^{-9}\text{ !}$
  Como $P' \in [10^{17}, 10^{18} - 1]$ já é estritamente menor que $10^{18}$, \textbf{qualquer} fator primo de $P'$ é automaticamente $< 10^{18}$! Logo, basta tomar $P' = P$ (os primeiros 18 dígitos exatos de $X$), fatorar $P'$ com \textbf{Pollard's Rho} em $\approx 2\text{ms}$ e somar $+E$ aos expoentes dos primos $2$ e $5$!}
  {Esta propriedade resolve a fatoração aproximada em uma única chamada de Pollard's Rho: como $P < 10^{18}$, todo fator primo de $P$ já é $< 10^{18}$, e $\frac{|X - P \cdot 10^E|}{X} < \frac{10^E}{10^{17} \cdot 10^E} = 10^{-17} \ll 10^{-9}$!}
  {Miller-Rabin: $\Ocomp{k \log^3 N}$ $\mid$ Pollard's Rho: $\Ocomp{N^{1/4}}$ ($\approx 1\text{ms}$ para $10^{18}$)}

\begin{lstlisting}
ull modmul64(ull a, ull b, ull m) { return (ull)((u128)a * b % m); }
ull modpow64(ull b, ull e, ull m) {
  ull r = 1;
  for (; e > 0; e >>= 1, b = modmul64(b, b, m)) if (e & 1) r = modmul64(r, b, m);
  return r;
}

/* 1. Miller-Rabin Deterministico para unsigned 64-bit (< 2^64) */
bool is_prime64(ull n) {
  if (n < 2 || n % 6 % 4 != 1) return (n | 1) == 3;
  static constexpr ull A[] = {2, 325, 9375, 28178, 450775, 9780504, 1795265022};
  ull s = countr_zero(n - 1), d = n >> s;
  for (ull a : A) {
    ull p = modpow64(a % n, d, n), i = s;
    while (p != 1 && p != n - 1 && a % n && i--) p = modmul64(p, p, n);
    if (p != n - 1 && i != s) return false;
  }
  return true;
}

/* 2. Algoritmo de Pollard's Rho (Brent + multiplicacao acumulada em blocos de 128) */
ull pollard_rho(ull n) {
  auto f = [n](ull x) { return modmul64(x, x, n) + 1; };
  ull x = 0, y = 0, t = 30, prd = 2, i = 1, q;
  while (t++ % 40 || std::gcd(prd, n) == 1) {
    if (x == y) x = ++i, y = f(x);
    if ((q = modmul64(prd, max(x, y) - min(x, y), n))) prd = q;
    x = f(x), y = f(f(y));
  }
  return std::gcd(prd, n);
}

void factorize64_rec(ull n, vector<ull>& fac) {
  if (n == 1) return;
  if (is_prime64(n)) { fac.pb(n); return; }
  ull d = pollard_rho(n);
  factorize64_rec(d, fac);
  factorize64_rec(n / d, fac);
}

/* Retorna lista ordenada de pares {primo, expoente} para n <= 10^18 */
vector<pair<ull, int>> factorize64(ull n) {
  vector<ull> fac;
  factorize64_rec(n, fac);
  sort(all(fac));
  vector<pair<ull, int>> res;
  for (ull p : fac) {
    if (res.empty() || res.back().fst != p) res.eb(p, 1);
    else res.back().snd++;
  }
  return res;
}

/* 3. Fatoracao aproximada via prefixo de 18 digitos para X de ate 1000 digitos em O(X_prefix^{1/4}) */
vector<pair<ull, int>> solve_fuzzy_factorization(const string& S_X) {
  int len = sz(S_X), take = min(len, 18), rem_zeros = len - take;
  ull pref_val = stoull(S_X.substr(0, take));
  map<ull, int> mp;
  for (auto [p, e] : factorize64(pref_val)) mp[p] += e;
  if (rem_zeros > 0) {
    mp[2] += rem_zeros;
    mp[5] += rem_zeros;
  }
  return vector<pair<ull, int>>(all(mp));
}
\end{lstlisting}
